\section{Определение абсолютно сходящегося ряда. Критерий Коши для рядов
комплексных
чисел.}\label{ux43eux43fux440ux435ux434ux435ux43bux435ux43dux438ux435-ux430ux431ux441ux43eux43bux44eux442ux43dux43e-ux441ux445ux43eux434ux44fux449ux435ux433ux43eux441ux44f-ux440ux44fux434ux430.-ux43aux440ux438ux442ux435ux440ux438ux439-ux43aux43eux448ux438-ux434ux43bux44f-ux440ux44fux434ux43eux432-ux43aux43eux43cux43fux43bux435ux43aux441ux43dux44bux445-ux447ux438ux441ux435ux43b.}

\textbf{Определение:} Говорят, что ряд
\(\sum\limits_{n=1}^{\infty}z_{n}\) сходится \emph{абсолютно}, если
\(\sum\limits_{n=1}^{\infty}|z_{n}|\) сходится.

\textbf{Теорема:} (Критерий Коши для рядов комплексных чисел) \[
\forall\varepsilon>0 \ \exists n_{\varepsilon}: \ \forall m,n\geq n_{\varepsilon}\implies |\underbrace{z_{m+1}+z_{m+2}+\dots+z_{n}}_{S_{n}-S_{m}}|<\varepsilon
\] это необходимое и достаточное условие сходимости ряда.
